%! Exponential Function Manipulation — Unit 4, Lesson 4.4 — Homework (Answer Key)
\documentclass[10pt]{article}
\usepackage{precalculus-article}
\usepackage{precalculus-key}   % pulls in -boxes; do NOT also load -boxes
\newcommand{\TallMath}[1]{$\displaystyle #1\rule[-1.4em]{0pt}{3.2em}$}

\begin{document}

\pageheader{Unit 4, Lesson 4.4}{Homework --- Answer Key}
\namedateperiod

\vspace{0.12in}

\begin{enumerate}[itemsep=18pt]

\item \textbf{Product Property: Horizontal Translation $\to$ Vertical Dilation.}
      Rewrite each expression in the form $a \cdot b^x$. Show work.

      \begin{enumerate}[label=(\alph*), itemsep=10pt]
        \item $f(x) = 2^{x+5}$

              \vspace{0.06in}
              $2^{x+5} = 2^x \cdot 2^{{\color{keyred}\mathbf{5}}} = $ \ans{32} $\cdot\; 2^x$

        \item $g(x) = 7^{x-2}$

              \vspace{0.06in}
              $7^{x-2} = 7^x \cdot 7^{{\color{keyred}\mathbf{-2}}} = $ \ans{$\dfrac{1}{49}$} $\cdot\; 7^x$
      \end{enumerate}

\item \textbf{Power Property: Horizontal Dilation $\to$ Base Change.}
      Rewrite each expression as a single base raised to the $x$ power. Show work.

      \begin{enumerate}[label=(\alph*), itemsep=10pt]
        \item $f(x) = 3^{4x}$

              \vspace{0.06in}
              $3^{4x} = (3^{{\color{keyred}\mathbf{4}}})^x = $ \ans{$81^x$}

        \item $g(x) = 25^{x/2}$

              \vspace{0.06in}
              $25^{x/2} = (25^{{\color{keyred}\mathbf{1/2}}})^x = $ \ans{$5^x$}
      \end{enumerate}

\item \textbf{Negative and Fractional Exponents.}
      Simplify each expression. Write without negative or fractional exponents where possible.

      \begin{enumerate}[label=(\alph*), itemsep=8pt]
        \item $5^{-2} = $ \ans{$\dfrac{1}{25}$}
        \item $27^{1/3} = $ \ans{3}
        \item $16^{3/4} = (\sqrt[4]{16})^3 = 2^3 = $ \ans{8}
      \end{enumerate}

\item \textbf{Show Two Functions Are Equivalent.}

      \vspace{0.06in}
      $f(x) = 4^{x+3/2} = 4^x \cdot 4^{{\color{keyred}\mathbf{3/2}}} = $ \ans{8} $\cdot\; 4^x$

      \quad ($4^{3/2} = (\sqrt{4})^3 = 2^3 = 8$)

      \vspace{0.06in}
      $g(x) = 8 \cdot 4^x \cdot \sqrt{4}^{\,-1} = 8 \cdot 4^x \cdot 2^{-1} = $ \ans{4} $\cdot\; 4^x$

      \begin{teachernote}
        Note: $f(x) = 8 \cdot 4^x$ and $g(x) = 4 \cdot 4^x$ are NOT equal --- check the problem
        statement. If $g(x) = 8 \cdot 4^x \cdot (\sqrt{4})^{-1} = 8 \cdot (1/2) \cdot 4^x = 4 \cdot 4^x$,
        these are not equal. The problem is intentionally asking students to check equality by
        computing both --- the answer is \textbf{No, they are not equivalent}. This tests careful
        algebraic work over blind assumption.
      \end{teachernote}

      Are they equivalent? \ans{No --- $f(x) = 8 \cdot 4^x \ne 4 \cdot 4^x = g(x)$}

\item \textbf{(AP-Style Multiple Choice)} Which expression is equivalent to $16^{x/4}$?

      Answer: \ans{(B)} \quad Reasoning: $16^{x/4} = (16^{1/4})^x = 2^x$.

      \begin{teachernote}
        $16^{1/4} = \sqrt[4]{16} = 2$ since $2^4 = 16$. Answer (C): $4^{x/2} = (4^{1/2})^x = 2^x$ is
        also equal --- but the intended answer matching the standard form $(b^c)^x \to b^x$ is (B).
        Accept (C) with full justification.
      \end{teachernote}

\item \textbf{Extension.} Find $b$ such that $f(x) = b^{2x}$ and $g(x) = 9 \cdot 3^x$ are the
      same function.

      \vspace{0.06in}
      Rewrite $f$: $b^{2x} = (b^2)^x$.
      For $f = g$, we need $(b^2)^x = 9 \cdot 3^x$.

      But $9 \cdot 3^x$ is not in the form $c^x$ unless we write $9 \cdot 3^x = 3^2 \cdot 3^x = 3^{x+2}$,
      which is not a single-base power of $x$.

      The functions can be equal only if $9 = 1$, which is impossible --- \textbf{there is no value of $b$}
      that makes $b^{2x}$ (form: $(b^2)^x$, no coefficient) equal to $9 \cdot 3^x$ (form: $a \cdot 3^x$ with $a \ne 1$).

      \ans{No solution.} \quad ($b^{2x}$ has $a = 1$; $g$ has $a = 9 \ne 1$.)

      \begin{teachernote}
        This problem tests whether students recognize that $b^{2x} = (b^2)^x$ always has $a = 1$
        (no vertical dilation), so it cannot equal $9 \cdot 3^x$. Stronger students may ask whether
        $b$ could be complex --- a rich discussion opportunity.
      \end{teachernote}

\end{enumerate}

\end{document}
